
\section{Document Data}

    \subsection{Front-Matter Footer Block*}

        The \inlinecode{\docinfo{}} command inserts a dedicated block at the bottom of the first column. Unlike floating elements, this block stays fixed in place — positioned visually like a footnote, although technically it is not one.
    
        The \inlinecode{\docinfo{}} block was created to give you greater control over supplementary document data, such as publication dates, licensing information, extended author details, and other front-matter notes. You can fully customize its content to suit your needs. If it’s not required, omit the command to remove this block entirely.
    
        \begin{note}
            This feature has been adapted to work seamlessly alongside the standard \verb|\thanks{}| command, which is commonly used to indicate corresponding authors, equal contributions, or other acknowledgments.
        \end{note}
    
        If \inlinecode{\thanks{}} is not used, \inlinecode{\docinfo{}} will automatically adjust its position. The implementation was designed to handle any combination of front-matter elements you might need when starting your document. In this example template, you’ll find a minimal demonstration that combines both to illustrate their interaction.

    \subsection{Headers and Footers}

        \subsubsection{Headers}

            On all pages except the first, the document title appear in the header. Since twoside layout is enabled, its position alternates depending on whether the page is odd or even.

        \subsubsection{Footers}

            On odd-numbered pages, five elements display supplementary information for your document:
    
            \begin{itemize}
                \item \inlinecode{\thepage} — Including the total number of pages (only for the first page).
                \item \inlinecode{\footinfo{}} — For a short title or custom note.
                \item \inlinecode{\theday{}} — Displays the current date.
                \item \inlinecode{\organization{}} — For a university, institute, company, or any institutional affiliation.
                \item \inlinecode{\leadauthor{}} — For the main author name et al.
            \end{itemize}

            On even-numbered pages, the page number appears alongside the content of \inlinecode{\footinfo{}} and, on the other side, the \inlinecode{\leadauthor{}}. 

            In contrast to the header, the footer elements do not alternate their positions. The footer is explicitly configured for either odd or even pages, resulting in a fixed layout.
            
            In all cases, you can freely rearrange or customize the order of these elements by modifying the class file to suit your needs. If any of them are not required, simply omit the corresponding command — the remaining elements will automatically adjust their spacing to remain evenly distributed.

